% Fragmento público generado desde propietarios íntegros.
% Residencia: 52.10.2.
% Fuente: ../incoming/radion_monografia_20260827/manuscrito/sections/11_modularidad_orbitas.tex:85-187
\subsubsection{Conservación areal en las órbitas de Kepler y Sommerfeld}

\paragraph{Equivalencia areal de las realizaciones orbital y oscilatoria}

La elipse de Kepler vive en espacio de configuración y responde a un potencial
central. La elipse del radión vive en espacio de fases y responde a una forma
de acción. Son objetos diferentes. Comparten una ley areal porque ambas
realizaciones poseen una dos-forma conservada.

Para una órbita kepleriana,
\begin{equation}
r(\phi)=\frac{p}{1+e\cos(\phi-\phi_0)},
\qquad
\frac{\diff A}{\diff t}=\frac{\ell}{2m}.
\label{eq:kepler}
\end{equation}
La tercera ley toma la forma
\[
T^2=\frac{4\pi^2\mu}{K}a^3.
\]
Para el oscilador en fase,
\begin{equation}
J=\oint p\,\diff q=\frac{2\pi E}{\omega}.
\label{eq:osc-action}
\end{equation}
Si la celda elemental es \(J=\hfull\), entonces
\begin{equation}
E=\hret\omega,\qquad \frac{J}{2\pi}=\hret.
\label{eq:E-hbarw}
\end{equation}
Ésta es la convergencia limpia entre acción de Planck, frecuencia angular y
radión.

Sommerfeld cuantiza integrales de acción sobre ciclos. La lectura HMT añade
la memoria de qué ciclo, hoja y retorno sobreviven. En una banda de Möbius,
un bucle visible puede portar \(\hfull/2\) y dos bucles restaurar \(\hfull\).
Ese factor topológico no se suma al índice de Maslov de una cuantización EBK;
pertenecen a correcciones con propietarios distintos.

\paragraph{Holonomía orbital y propagaciones avanzada y retardada}

Una órbita que acumula holonomía satisface una condición del tipo
\begin{equation}
\epsilon_\gamma\exp\left(\frac{iJ_\gamma}{\hret}\right)=1.
\label{eq:holonomy-action}
\end{equation}
El signo \(\epsilon_\gamma\) conserva la hoja; \(J_\gamma\) conserva la
acción. El adelanto o retardo apsidal se interpreta como desacoplamiento entre
el cierre angular visible y el retorno completo. El radión suministra la
unidad orientada con la que se mide esa diferencia.

\subsubsection{Composición de amplitudes de Schrödinger y suma de caminos}

\paragraph{Cuatro realizaciones de una misma fase}

Una vez fijado \(\hret\), distintas ecuaciones convencionales reconocen
aspectos del mismo transporte:
\begin{enumerate}
\item Schrödinger publica la fase mediante
\(i\hret\partial_t\psi=H\psi\).
\item Pauli añade la doble hoja de espín y su acoplamiento magnético.
\item Dirac linealiza la relación energía--momento y organiza partícula y
antipartícula en una representación causal.
\item La suma de caminos asigna a cada historia la fase
\(\exp(iS[\gamma]/\hret)\).
\end{enumerate}
Estas teorías no generan retrospectivamente \(\hret\). Reciben la unidad de
acción y realizan diferentes contenidos del estado enriquecido.

\paragraph{Funcional de fase sobre el espacio de caminos}

En la formulación de Feynman,
\begin{equation}
\mathcal K(b,a)=\int_{\gamma:a\to b}
\exp\left(\frac{i}{\hret}S[\gamma]\right)\mathcal D\gamma.
\label{eq:path-integral}
\end{equation}
Un incremento de acción \(\hret\) cambia la fase en un radián. Un incremento
\(\hfull\) la cambia en una vuelta. El teorema
\(\mathcal A(\theta)=\hret\theta\) convierte esta relación en geometría de
área: la fase de un camino puede leerse como acción barrida en unidades de
radión.

La inversión \(C_M=-\operatorname{rev}\) sobre palabras dodecafásicas y la
identidad \eqref{eq:time-reversal} organizan rutas conjugadas. En una
realización de absorbedor se separan
\[
G_{\mathrm{sym}}=\tfrac12(G_{\mathrm{ret}}+G_{\mathrm{adv}}),
\qquad
G_{\mathrm{odd}}=G_{\mathrm{ret}}-G_{\mathrm{adv}}.
\]
El primer sector es par bajo intercambio retardado--avanzado; el segundo es
impar. Esto reproduce la separación entre energía/forma par y acción
orientada impar del radión cuando el transductor preserva frontera y forma
simpléctica.

\begin{statusbox}
La teoría clásica del absorbedor y la interpretación Feynman--Stückelberg de
antipartículas son realizaciones históricas diferentes. La primera usa una
interacción temporalmente simétrica; la segunda reorganiza propagación en la
teoría relativista. HMT aporta un antecedente común de reversión, pero no las
funde sin un mapa de propagadores.
\end{statusbox}
